\documentclass[12pt,a4paper]{ctexart}

\usepackage{geometry}
\usepackage{fancyhdr}

% 页面设置：A4，上下2.54cm，左右3.175cm，匹配原Word
\geometry{top=2.54cm, bottom=2.54cm, left=3.175cm, right=3.175cm, headheight=15pt}

% 辅助命令：下划线空白
\newcommand{\blankline}[1]{\underline{\hspace{#1}}}

% 页眉设置
\pagestyle{fancy}
\fancyhf{}
\fancyhead[C]{齐鲁工业大学（山东省科学院）数学与人工智能学部}
\renewcommand{\headrulewidth}{0.4pt}

\begin{document}

% ====== 标题行 ======
\noindent
\textbf{齐鲁工业大学（山东省科学院）实验报告}\hfill 成绩\blankline{3cm}

\vspace{0.6cm}

% ====== 基本信息行1 ======
\noindent
课程名称\blankline{2.5cm}
指导教师\blankline{2.5cm}
实验日期\blankline{2.5cm}

\vspace{0.5cm}

% ====== 基本信息行2 ======
\noindent
学部（学院）\blankline{3cm}
专业班级\blankline{2.5cm}
实验地点\blankline{2.5cm}

\vspace{0.5cm}

% ====== 基本信息行3 ======
\noindent
学生姓名\blankline{2.5cm}
学号\blankline{2.5cm}
同组人\blankline{3cm}

\vspace{0.6cm}

% ====== 实验项目名称 ======
\noindent
实验项目名称\blankline{8cm}

\vspace{0.8cm}

% ====== 各章节区域 ======

\noindent
\textbf{一、实验目的和要求}

\vspace{2cm}

\noindent
\textbf{二、实验原理}

\vspace{2cm}

\noindent
\textbf{三、主要仪器设备}

\vspace{2cm}

\noindent
\textbf{四、实验方法与步骤}

\vspace{2cm}

\noindent
\textbf{五、实验内容及实现源代码}

\vspace{2cm}

\noindent
\textbf{六、实验数据记录、处理及结果分析}

\vspace{2cm}

\noindent
\textbf{讨论、心得}

\vspace{2cm}

\end{document}
